\documentclass[11pt,a4paper]{article}
\usepackage{fontspec,amsmath,amssymb,amsthm,microtype}
\setmainfont{lmroman10-regular.otf}[BoldFont=lmroman10-bold.otf,ItalicFont=lmroman10-italic.otf,BoldItalicFont=lmroman10-bolditalic.otf]
\usepackage{xcolor}
\usepackage[most]{tcolorbox}
\usepackage{mathrsfs}
\providecommand{\lt}{<}
\providecommand{\gt}{>}
\setmonofont{DejaVuSansMono.ttf}[BoldFont=DejaVuSansMono-Bold.ttf,ItalicFont=DejaVuSansMono-Oblique.ttf,BoldItalicFont=DejaVuSansMono-BoldOblique.ttf,Scale=0.86]
\newfontfamily\leanfallback{STIXTwoMath-Regular.otf}[Scale=0.86]
\definecolor{leanrule}{gray}{0.62}
\definecolor{leanback}{gray}{0.962}
\definecolor{leantitle}{gray}{0.22}
\definecolor{leanmuted}{gray}{0.38}
\newtcolorbox{leanbox}[1]{enhanced,breakable,colback=leanback,colframe=leanrule,
  boxrule=0.4pt,arc=0.6mm,left=2.4mm,right=2.4mm,top=1.8mm,bottom=1.8mm,
  before skip=0.8em,after skip=1em,fonttitle=\small\bfseries,coltitle=leantitle,
  colbacktitle=leanback,titlerule=0.3pt,title={#1}}
\newlength{\leancharwidth}
\newenvironment{leancode}{\par\smallskip\begingroup\ttfamily\footnotesize\raggedright
  \settowidth{\leancharwidth}{M}\setlength{\parskip}{0pt}\setlength{\parindent}{0pt}
  \frenchspacing}{\par\endgroup\smallskip}
\newcommand{\leanline}[2]{\hangindent=\dimexpr#1\leancharwidth+4\leancharwidth\relax
  \hangafter=1\noindent\hspace*{#1\leancharwidth}#2\par}
\newcommand{\leansource}[1]{\par{\footnotesize\color{leanmuted}#1}\par}
\newcommand{\leanlabel}[1]{\par\smallskip\noindent{\small\bfseries #1}\par}
\usepackage[margin=25mm]{geometry}
\usepackage[unicode,hidelinks]{hyperref}
\hypersetup{pdfauthor={LangevinLean Project}}
\setlength{\emergencystretch}{3em}
\setlength{\parskip}{0.4em}
\setlength{\parindent}{0pt}
\allowdisplaybreaks[2]
\raggedbottom
\title{\texorpdfstring{Smoothing by arbitrary independent centered coordinates}{Smoothing by arbitrary independent centered coordinates}}
\author{LangevinLean Project}
\date{}
\begin{document}\maketitle
Draft mathematical exposition; these notes may contain mistakes.


This module proves the full expectation conclusion of the manuscript smoothing lemma. The coordinates may have any distributions with zero mean and unit second moment; Gaussianity is used only in the final specialization. The bound uses the exact coordinate Taylor suprema and their arithmetic average. All potential, score, and product integrability follows from the supplied coordinate second moments and the gradient Lipschitz assumption.



Earlier results used here: \cite[Section~3.13]{LangevinLean2026ULASC}; \cite[Section~3.4]{LangevinLean2026ULASC}.

\section{\texorpdfstring{Build square-integrable partial coordinate sums}{Build square-integrable partial coordinate sums}}


For a finite set \(S\) of coordinate indices, define \(Z_S\) as the sum of \(\xi_j\) \(e_j\) over \(j\) in \(S\). Each scalar coordinate is square integrable, and the map \(a\) to \(ae_j\) is continuous linear, so every vector summand and every finite partial sum is square integrable. The actual gradient evaluated at an \(L^{2}\) state is also \(L^{2}\): apply the existing Lipschitz drift lift to the canonical \(L^{2}\) representative, then transfer back along its almost-everywhere equality. Adding the deterministic center and scaling preserve these moments.

\begin{leanbox}{Lean: \texttt{memLp\_\allowbreak{}gradient\_\allowbreak{}comp}, \texttt{coordinatePartial}, \texttt{memLp\_\allowbreak{}coordinatePartial}}
\leansource{\texttt{memLp\_\allowbreak{}gradient\_\allowbreak{}comp} \quad AverageLMC/\allowbreak{}Smoothing.lean, lines 28–35}
\begin{leancode}
\leanline{0}{private\ theorem\ memLp\_gradient\_comp\ \{p\ :\ ℕ\}\ (μ\ :\ Measure\ Ω)\ [IsProbabilityMeasure\ μ]}
\leanline{4}{(f\ :\ State\ p\ →\ ℝ)\ (M\ :\ ℝ≥0)\ (hLip\ :\ LipschitzWith\ M\ (gradient\ f))}
\leanline{4}{(X\ :\ Ω\ →\ State\ p)\ (hX\ :\ MemLp\ X\ 2\ μ)\ :}
\leanline{4}{MemLp\ (fun\ ω\ ↦\ gradient\ f\ (X\ ω))\ 2\ μ}
\end{leancode}
\leansource{\texttt{coordinatePartial} \quad AverageLMC/\allowbreak{}Smoothing.lean, lines 37–39}
\begin{leancode}
\leanline{0}{private\ def\ coordinatePartial\ \{p\ :\ ℕ\}\ (ξ\ :\ Fin\ p\ →\ Ω\ →\ ℝ)}
\leanline{4}{(s\ :\ Finset\ (Fin\ p))\ (ω\ :\ Ω)\ :\ State\ p\ :=}
\leanline{2}{∑\ j\ ∈\ s,\ ξ\ j\ ω\ •\ coordinateVector\ j}
\end{leancode}
\leansource{\texttt{memLp\_\allowbreak{}coordinatePartial} \quad AverageLMC/\allowbreak{}Smoothing.lean, lines 41–46}
\begin{leancode}
\leanline{0}{private\ theorem\ memLp\_coordinatePartial\ \{p\ :\ ℕ\}\ (μ\ :\ Measure\ Ω)}
\leanline{4}{(ξ\ :\ Fin\ p\ →\ Ω\ →\ ℝ)\ (hξ\ :\ ∀\ j,\ MemLp\ (ξ\ j)\ 2\ μ)\ (s\ :\ Finset\ (Fin\ p))\ :}
\leanline{4}{MemLp\ (coordinatePartial\ ξ\ s)\ 2\ μ}
\end{leancode}
\leanlabel{How the Lean proof follows the argument}
\texttt{mem\allowbreak{}Lp\_\allowbreak{}gradient\_\allowbreak{}comp} gives the formal counterpart of this step. The main intermediate facts are \texttt{h} in that order. Named proof applications include \texttt{h.\allowbreak{}ae\_\allowbreak{}eq}; the adjacent proof shows their hypotheses and order of use.




\texttt{coordinate\allowbreak{}Partial} names the mathematical object defined above. Its defining expression appears in the adjacent statement; subsequent identities may unfold this definition directly.




\texttt{mem\allowbreak{}Lp\_\allowbreak{}coordinate\allowbreak{}Partial} gives the formal counterpart of this step. Named proof applications include \texttt{mem\allowbreak{}Lp\_\allowbreak{}finset\allowbreak{}Sum}; the adjacent proof shows their hypotheses and order of use.
\end{leanbox}


Earlier results used here: \cite[Section~3.13.4]{LangevinLean2026ULASC}.

\section{\texorpdfstring{Insert one independent coordinate and eliminate its linear term}{Insert one independent coordinate and eliminate its linear term}}


Induct on the finite coordinate set. When inserting \(i\) outside \(S\), apply the exact coordinate model at \(x+c\) \(Z_S\) with displacement \(c\) \(\xi_i\). It gives

\[f(x+cZ_{S\cup\{i\}})\le f(x+cZ_S)+c\xi_i\partial_i f(x+cZ_S)+\tfrac12 M_i c^2\xi_i^2.\]

Independence of the scalar coordinates survives the maps \(a\) to \(ae_j\). Consequently their partial vector sum is independent of the new coordinate. Composing with the actual gradient coordinate proves independence of \(\xi_i\) and the displayed derivative term. Their product is integrable by the \(L^{2}\) Holder inequality, and its integral factors; the zero first moment makes it vanish. The quadratic term integrates to \(M_i\) \(c^2/2\). Adding this to the inductive bound proves the finite-set estimate, with no assumed stochastic error inequality.




\textbf{Induction in detail.} At the empty set the perturbation is zero and the expected potential is \(f(x)\). Suppose the bound holds for \(S\), and choose \(i\notin S\). Set \(G_i=\partial_i f(x+cZ_S)\). Global gradient Lipschitz continuity and the finite second moment of \(Z_S\) imply \(G_i\in L^2\); hence Cauchy–Schwarz makes \(\xi_iG_i\) integrable. The variable \(G_i\) is a measurable function of the coordinates in \(S\), so independence gives \(\mathbb E[\xi_iG_i]=(\mathbb E\xi_i)(\mathbb EG_i)=0\). The coordinate model then gives the previous expectation plus \(M_ic^2/2\), since \(\mathbb E\xi_i^2=1\). The induction hypothesis and \(\sum_{j\in S\cup\{i\}}M_j=\sum_{j\in S}M_j+M_i\) complete the step. This holds for every real \(c\), including \(c=0\), and for arbitrary independent centered unit-second-moment coordinates with measurable representatives. Gaussianity is used only in the later specialization. \(\square\)

\begin{leanbox}{Lean: \texttt{integral\_\allowbreak{}coordinatePartial\_\allowbreak{}le}}
\leansource{\texttt{integral\_\allowbreak{}coordinatePartial\_\allowbreak{}le} \quad AverageLMC/\allowbreak{}Smoothing.lean, lines 50–117}
\begin{leancode}
\leanline{0}{theorem\ integral\_coordinatePartial\_le\ \{p\ :\ ℕ\}\ (μ\ :\ Measure\ Ω)\ [IsProbabilityMeasure\ μ]}
\leanline{4}{(f\ :\ State\ p\ →\ ℝ)\ (M\ :\ ℝ≥0)\ (hf\ :\ Differentiable\ ℝ\ f)}
\leanline{4}{(hLip\ :\ LipschitzWith\ M\ (gradient\ f))}
\leanline{4}{(ξ\ :\ Fin\ p\ →\ Ω\ →\ ℝ)\ (hξm\ :\ ∀\ j,\ Measurable\ (ξ\ j))}
\leanline{4}{(hξ\ :\ ∀\ j,\ MemLp\ (ξ\ j)\ 2\ μ)\ (hind\ :\ iIndepFun\ ξ\ μ)}
\leanline{4}{(hmean\ :\ ∀\ j,\ ∫\ ω,\ ξ\ j\ ω\ ∂μ\ =\ 0)}
\leanline{4}{(hvar\ :\ ∀\ j,\ ∫\ ω,\ ξ\ j\ ω\ \textasciicircum{}\ 2\ ∂μ\ =\ 1)}
\leanline{4}{(x\ :\ State\ p)\ (c\ :\ ℝ)\ (s\ :\ Finset\ (Fin\ p))\ :}
\leanline{4}{(∫\ ω,\ f\ (x\ +\ c\ •\ coordinatePartial\ ξ\ s\ ω)\ ∂μ)\ ≤}
\leanline{6}{f\ x\ +\ c\ \textasciicircum{}\ 2\ /\ 2\ *\ ∑\ j\ ∈\ s,\ coordinateSmoothness\ f\ j}
\end{leancode}
\leanlabel{How the Lean proof follows the argument}
\texttt{integral\_\allowbreak{}coordinate\allowbreak{}Partial\_\allowbreak{}le} derives the partial-state moments \texttt{h\allowbreak{}X} and potential integrability \texttt{h\allowbreak{}F} before the finite-set induction. In the insert branch \texttt{G} is the new score coordinate and \texttt{h\allowbreak{}G} proves its square integrability. \texttt{h\allowbreak{}IG} obtains its independence from the new coordinate by composing independence of the finite partial sum with the score function. \texttt{hprod} supplies product integrability, and \texttt{hzero} factors the integral and uses centering. The coordinate Taylor model is then integrated, its squared-coordinate integral is replaced by the unit-moment hypothesis, and the induction hypothesis supplies the sum over the old coordinates.
\end{leanbox}

\section{\texorpdfstring{Sum every coordinate and recover the exact average-smoothness coefficient}{Sum every coordinate and recover the exact average-smoothness coefficient}}


For the full index set, the coordinate sum is exactly the Euclidean vector with coordinates \(\xi_j\). Choose \(c\) equal to the square root of \(2s\), whose square is \(2s\) because \(s\) is nonnegative. Smooth potentials have quadratic growth, so the perturbed potential and its difference from the deterministic center are integrable. The finite-set theorem yields

\[\mathbb E[f(x+\sqrt{2s}\,\xi)-f(x)]\le s\sum_j M_j=M_{\rm av}ps.\]

The final equality uses the literal definition \(M_{\mathrm{av}}=p^{-1}\sum_jM_j\) and the positive dimension.

\begin{leanbox}{Lean: \texttt{integral\_\allowbreak{}coordinate\_\allowbreak{}smoothing}}
\leansource{\texttt{integral\_\allowbreak{}coordinate\_\allowbreak{}smoothing} \quad AverageLMC/\allowbreak{}Smoothing.lean, lines 121–151}
\begin{leancode}
\leanline{0}{theorem\ integral\_coordinate\_smoothing\ \{p\ :\ ℕ\}\ (hp\ :\ 0\ <\ p)}
\leanline{4}{(μ\ :\ Measure\ Ω)\ [IsProbabilityMeasure\ μ]}
\leanline{4}{(f\ :\ State\ p\ →\ ℝ)\ (M\ :\ ℝ≥0)\ (hf\ :\ Differentiable\ ℝ\ f)}
\leanline{4}{(hLip\ :\ LipschitzWith\ M\ (gradient\ f))}
\leanline{4}{(ξ\ :\ Fin\ p\ →\ Ω\ →\ ℝ)\ (hξm\ :\ ∀\ j,\ Measurable\ (ξ\ j))}
\leanline{4}{(hξ\ :\ ∀\ j,\ MemLp\ (ξ\ j)\ 2\ μ)\ (hind\ :\ iIndepFun\ ξ\ μ)}
\leanline{4}{(hmean\ :\ ∀\ j,\ ∫\ ω,\ ξ\ j\ ω\ ∂μ\ =\ 0)}
\leanline{4}{(hvar\ :\ ∀\ j,\ ∫\ ω,\ ξ\ j\ ω\ \textasciicircum{}\ 2\ ∂μ\ =\ 1)}
\leanline{4}{(x\ :\ State\ p)\ (s\ :\ ℝ)\ (hs\ :\ 0\ ≤\ s)\ :}
\leanline{4}{(∫\ ω,\ f\ (x\ +\ Real.sqrt\ (2\ *\ s)\ •\ WithLp.toLp\ 2\ (fun\ j\ ↦\ ξ\ j\ ω))\ -\ f\ x\ ∂μ)\ ≤}
\leanline{6}{averageSmoothness\ f\ *\ (p\ :\ ℝ)\ *\ s}
\end{leancode}
\leanlabel{How the Lean proof follows the argument}
\texttt{integral\_\allowbreak{}coordinate\_\allowbreak{}smoothing} gives the formal counterpart of this step. The main intermediate facts are \texttt{hi}, \texttt{heq}, \texttt{h\allowbreak{}X}, \texttt{h\allowbreak{}F}, \texttt{hsum} in that order. Its remaining equalities are discharged by the simplifications displayed in the proof.
\end{leanbox}

\section{\texorpdfstring{Specialize to the actual standard Gaussian measure}{Specialize to the actual standard Gaussian measure}}


Take the product of \(n=p\) real Gaussian laws with mean zero and variance one. Each coordinate evaluation preserves its marginal Gaussian measure; the product coordinates are mutually independent, square integrable, centered, and have second moment one. Apply the arbitrary-coordinate theorem on this product. Its pushforward is identified by the Euclidean coordinate map with the actual standard Gaussian law, so the change-of-measure integral is precisely the Gaussian smoothing integral used by the LMC transition.

\begin{leanbox}{Lean: \texttt{integral\_\allowbreak{}stdGaussian\_\allowbreak{}smoothing}}
\leansource{\texttt{integral\_\allowbreak{}stdGaussian\_\allowbreak{}smoothing} \quad AverageLMC/\allowbreak{}Smoothing.lean, lines 154–175}
\begin{leancode}
\leanline{0}{theorem\ integral\_stdGaussian\_smoothing\ \{p\ :\ ℕ\}\ (hp\ :\ 0\ <\ p)}
\leanline{4}{(f\ :\ State\ p\ →\ ℝ)\ (M\ :\ ℝ≥0)\ (hf\ :\ Differentiable\ ℝ\ f)}
\leanline{4}{(hLip\ :\ LipschitzWith\ M\ (gradient\ f))\ (x\ :\ State\ p)\ (s\ :\ ℝ)\ (hs\ :\ 0\ ≤\ s)\ :}
\leanline{4}{(∫\ z,\ f\ (x\ +\ Real.sqrt\ (2\ *\ s)\ •\ z)\ -\ f\ x\ ∂stdGaussian\ (State\ p))\ ≤}
\leanline{6}{averageSmoothness\ f\ *\ (p\ :\ ℝ)\ *\ s}
\end{leancode}
\leanlabel{How the Lean proof follows the argument}
\texttt{integral\_\allowbreak{}std\allowbreak{}Gaussian\_\allowbreak{}smoothing} gives the formal counterpart of this step. The main intermediate facts are \texttt{\ensuremath{\mu}}, \texttt{hl}, \texttt{h\ensuremath{\xi}}, \texttt{hmean}, \texttt{hvar}, \texttt{h} in that order.
\end{leanbox}


Earlier results used here: \cite[Section~3.4.1]{LangevinLean2026ULASC}.

\begin{thebibliography}{ULA-LSI}
\bibitem[ULA-SC]{LangevinLean2026ULASC}
LangevinLean Project. \emph{Unadjusted Langevin under Strong Convexity}. LL-01, 2026. Draft report, content revision d50c0f64815b. \url{https://odakyildiz.github.io/LangevinLean/docs/pdf/UnadjustedLangevin/FullProof.pdf}.
\end{thebibliography}

\end{document}
